\documentclass[UTF8,12pt]{ctexart}
\usepackage[a4paper,margin=24mm]{geometry}
\usepackage{amsmath,amssymb,bm,graphicx,adjustbox}
\newcommand{\fitmath}[1]{\begin{adjustbox}{max width=\linewidth}$\displaystyle #1$\end{adjustbox}}
\setlength{\parindent}{0pt}
\setlength{\parskip}{.7em}
\sloppy
\allowdisplaybreaks
\begin{document}
\section*{2018 年考研数学三 · 参考答案}
\subsection*{原 PDF 第 10 页}

\subsection*{2018年真题参考答案}

\subsection*{一、选择题}

（1）D。 （2）D。 （3）C。 （4）D。 （5）A。 （6）A。 （7）A。 （8）B。


\subsection*{二、填空题}

（9）\(\displaystyle y=4x-3\)。 （10）\(\displaystyle e^x\arcsin\sqrt{1-e^{2x}}-\sqrt{1-e^{2x}}+C\)。 （11）\(\displaystyle y_x=C2^x-5\)。 （12）\(\displaystyle 2e\)。 （13）\(\displaystyle 2\)。 （14）\(\displaystyle \dfrac{\displaystyle 1}{\displaystyle 3}\)。


\subsection*{三、解答题}

（15）\(\displaystyle a=1,\allowbreak b=1\)。 （16）\(\displaystyle \dfrac{\displaystyle \sqrt3}{\displaystyle 32}(\pi-2)\)。


（17）三个图形的面积之和存在最小值，最小值为 \(\displaystyle \dfrac{\displaystyle 1}{\displaystyle \pi+4+3\sqrt3}\)。


（18）\(\displaystyle a_n=\begin{cases}\dfrac{\displaystyle (-1)^k4^k}{\displaystyle (2k)!}-2k-1,\allowbreak &n=2k,\allowbreak \\2k+2,\allowbreak &n=2k+1,\allowbreak \end{cases}\) 其中 \(\displaystyle k\) 为非负整数。


（19）证明略。\(\displaystyle \lim_{n\to\infty}x_n=0\)。


（20）（I）当 \(\displaystyle a\ne2\) 时，\(\displaystyle f(x_1,\allowbreak x_2,\allowbreak x_3)=0\) 的解为 \(\displaystyle (x_1,\allowbreak x_2,\allowbreak x_3)^{\mathrm T}=(0,\allowbreak 0,\allowbreak 0)^{\mathrm T}\)；当 \(\displaystyle a=2\) 时，解为 \(\displaystyle (x_1,\allowbreak x_2,\allowbreak x_3)^{\mathrm T}=k(-2,\allowbreak -1,\allowbreak 1)^{\mathrm T}\)，其中 \(\displaystyle k\) 为任意常数。（II）当 \(\displaystyle a\ne2\) 时，\(\displaystyle f=y_1^2+y_2^2+y_3^2\)；当 \(\displaystyle a=2\) 时，\(\displaystyle f=z_1^2+z_2^2\)。


（21）（I）\(\displaystyle a=2\)。（II）满足 \(\displaystyle AP=B\) 的可逆矩阵为 \(\displaystyle P=\begin{pmatrix}-6k_1+3&-6k_2+4&-6k_3+4\\2k_1-1&2k_2-1&2k_3-1\\k_1&k_2&k_3\end{pmatrix}\)，其中 \(\displaystyle k_1,\allowbreak k_2,\allowbreak k_3\) 为任意常数，且 \(\displaystyle k_2\ne k_3\)。


（22）（I）\(\displaystyle \operatorname{Cov}(X,\allowbreak Z)=\lambda\)。（II）\(\displaystyle P\{Z=i\}=\begin{cases}\dfrac{\displaystyle 1}{\displaystyle 2}\dfrac{\displaystyle \lambda^i e^{-\lambda}}{\displaystyle i!},\allowbreak &i>0,\allowbreak \\e^{-\lambda},\allowbreak &i=0,\allowbreak \\\dfrac{\displaystyle 1}{\displaystyle 2}\dfrac{\displaystyle \lambda^{-i}e^{-\lambda}}{\displaystyle (-i)!},\allowbreak &i<0.\end{cases}\)


（23）（I）\(\displaystyle \hat\sigma=\dfrac{\displaystyle 1}{\displaystyle n}\displaystyle \sum_{i=1}^n|X_i|\)。（II）\(\displaystyle E(\hat\sigma)=\sigma\)，\(\displaystyle D(\hat\sigma)=\dfrac{\displaystyle \sigma^2}{\displaystyle n}\)。


\clearpage
\end{document}
